\section{The same sample exponent without clipping}
\label{sec:raw-gaussian-oracle}

The Gaussian-prior model and likelihood family in Section~\ref{sec:gaussian-oracle}, with the calibration below, also admit
an oracle inequality for its \emph{raw} real-valued predictions.  This
covers both posterior-averaged risk and the ordinary posterior mean
$\int z_w\,d\Pi_{\sigma,N,S}$, with no clipping in either the likelihood
or the predictor.  The improvement uses a standard uniform bound for
smooth losses.  It has the same sample exponent as the elementary
clipped argument, at the cost of additional logarithms.

We retain the architecture, energy $Q_N$, prior, likelihood, and risks
from Section~\ref{sec:gaussian-oracle}.  There are $H$ hidden layers and
one linear readout.  All labels lie in $[-1,1]$, and the joint law of
inputs and labels is otherwise arbitrary.  A teacher $g$ of width $m$
and weights bounded by $B$ has clones of energy at most
$E=d(d+1)B^2+2Hm^2B^2$ at every learner width $N\geq m$, and satisfies
$|g|\leq U=mB$.

\subsection{The smooth-loss inequality}

We use the following form of Theorem~1 of
Srebro, Sridharan, and Tewari~\cite[Theorem~1, p.~4 of the full-length version]{sst}.  If a nonnegative loss is
$s$-smooth in its real-valued prediction and is bounded by $b$ on a
function class $\mathcal F$, then, with probability at least $1-\eta$,
simultaneously for $f\in\mathcal F$,
\begin{align}
 R(f)-\widehat R_S(f)
 &\leq C_0\biggl[
  \sqrt{\widehat R_S(f)}
    \left(\sqrt{s}\log^{3/2}(n)\mathfrak R_n(\mathcal F)
           +\sqrt{\frac{b\log(1/\eta)}n}\right)
 \notag\\[-2pt]
 &\hspace{37mm}
  +s\log^3(n)\mathfrak R_n(\mathcal F)^2
  +\frac{b\log(1/\eta)}n\biggr].
 \label{eq:raw-sst}
\end{align}
Here $\mathfrak R_n$ is the worst-case Rademacher complexity and $C_0$ is
universal; $C_0=10^5$ is sufficient.  This result assumes neither a convex
hypothesis class nor a realizable data law.  In particular, it does not
assume a Bernstein condition on excess loss relative to the teacher.

The raw square loss is nonnegative and $2$-smooth.  On $\mathcal Z_K$
from~\eqref{eq:gaussian-Dn}, its range is bounded by
$(\sqrt K+1)^2$, since $|z_w|\leq\sqrt{Q_N(w)}$.
The class $\mathcal Z_K$ is symmetric under readout sign reversal, so
its Rademacher complexity with an absolute value inside the supremum,
as used in~\cite{sst}, equals the convention used in
Lemma~\ref{lem:gaussian-capacity}.

\subsection{The theorem for two or more hidden layers}

For $n\geq2$, $0<\delta<1$, and $H\geq2$, define
\begin{gather}
 p=H/2,\qquad J_n=\lceil\log_2 n\rceil,\qquad
 t_n=\log\frac{2(J_n+1)}\delta,\qquad q_n=t_n/n,\notag\\
 A_n=\frac{\sqrt2\,D_n\log^{3/2}(2n)}{H^p\sqrt n},\qquad
 a_n=2^p A_n,\notag\\
 r_n=C_0\bigl(A_n+2\sqrt{q_n}+A_n^2+4q_n\bigr).
 \label{eq:raw-constants}
\end{gather}
For a fixed teacher $g$, put
\begin{equation}
 b_n(g,\delta)
  =(U+1)\sqrt{\frac{2R(g)\log(2/\delta)}n}
     +\frac{2(U+1)^2\log(2/\delta)}{3n}.
 \label{eq:raw-teacher-concentration}
\end{equation}
The slightly conservative coefficient $2/3$ makes this a direct
consequence of the usual Bernstein tail inequality.

\begin{theorem}[A raw-output oracle inequality at every learner width]
\label{thm:raw-gaussian-oracle}
Let $H\geq2$, and suppose the likelihood parameter $\lambda$ satisfies
\begin{equation}
 \frac1n\leq\lambda\leq1,\qquad
 \lambda^p\geq8C_0a_n,\qquad \lambda\geq32C_0q_n.
 \label{eq:raw-calibration}
\end{equation}
For every fixed teacher $g$ as above, with probability at least
$1-\delta$ over the sample, simultaneously for every $N\geq m$,
\begin{equation}
 \limsup_{\sigma\downarrow0}
   \int R(z_w)\,\Pi_{\sigma,N,S}(dw)
 \leq R(g)+\lambda E+r_n+\frac{8C_0^2q_n}{\lambda}
                 +b_n(g,\delta).
 \label{eq:raw-oracle}
\end{equation}
The same right-hand side bounds the limiting risk of the ordinary
posterior mean $\int z_w\,d\Pi_{\sigma,N,S}$.  It also bounds the limiting
risk after clipping that mean, and the limiting risk of the posterior
mean of clipped outputs.

An explicit choice is
\begin{equation}
 \lambda_n=\max\bigl\{(8C_0a_n)^{1/p},\ 32C_0q_n\bigr\},
 \label{eq:raw-explicit-schedule}
\end{equation}
for all sufficiently large $n$ that $\lambda_n\leq1$.  At fixed
$d,H,\delta$, this gives an excess-risk bound of order
$\widetilde O(n^{-1/H})$ for every fixed teacher.  More explicitly,
\[
 \lambda_n
  =O_H\!\left[
      \left(\frac{D_n^2\log^3(2n)}n\right)^{1/H}
      +\frac{\log(\log(2n)/\delta)}n\right].
\]
The tilde suppresses logarithms; it does not suppress a learner-width
factor.  There is no dependence on $N$ in the displayed bound.
\end{theorem}

\begin{proof}
Apply~\eqref{eq:raw-sst} to each deterministic energy class
$\mathcal Z_{2^j}$, $0\leq j\leq J_n$, using failure probability
$\delta/[2(J_n+1)]$.  A union bound gives one event of probability at
least $1-\delta/2$ for all these classes and all learner widths.

Let $w$ minimize $\widehat R_S(z_w)+\lambda Q_N(w)$.  Comparing with
the zero network gives
\begin{equation}
 \widehat R_S(z_w)+\lambda Q_N(w)
   \leq\widehat R_S(0)\leq1.
 \label{eq:raw-training-budget}
\end{equation}
Write $k=Q_N(w)$.  If $k\leq1$, the event for $\mathcal Z_1$ implies
$R(z_w)-\widehat R_S(z_w)\leq r_n$.
If $k>1$, choose $j$ with $2^{j-1}<k\leq2^j$.  This is possible because
$k\leq1/\lambda\leq n$.  The shell radius $K=2^j$ satisfies $K\leq2k$.
The loss bound obeys $(\sqrt K+1)^2\leq4K\leq8k$.
Substituting Lemma~\ref{lem:gaussian-capacity} and
$\widehat R_S(z_w)\leq1$ into~\eqref{eq:raw-sst} yields
\begin{equation}
 R(z_w)-\widehat R_S(z_w)
 \leq C_0\left[a_n k^p+a_n^2 k^{2p}
                    +2\sqrt2\sqrt{kq_n}+8kq_n\right].
 \label{eq:raw-shell-bound}
\end{equation}
Since $p\geq1$, $k\leq1/\lambda$, and
$\lambda^p\geq8C_0a_n$, the first two terms satisfy
\[
 C_0a_n k^p\leq\frac{\lambda k}8,
 \qquad
 C_0a_n^2 k^{2p}\leq\frac{\lambda k}{64C_0}
                     \leq\frac{\lambda k}{64}.
\]
The remaining terms obey
\[
 8C_0kq_n\leq\frac{\lambda k}4,
 \qquad
 2\sqrt2 C_0\sqrt{kq_n}
  \leq\frac{\lambda k}4+\frac{8C_0^2q_n}{\lambda}.
\]
Consequently both the $k\leq1$ and $k>1$ cases imply
\begin{equation}
 R(z_w)\leq\widehat R_S(z_w)+\lambda Q_N(w)
                      +r_n+\frac{8C_0^2q_n}{\lambda}.
 \label{eq:raw-absorbed-deviation}
\end{equation}

Compare the penalized empirical objective with the exact teacher clone:
\[
 \widehat R_S(z_w)+\lambda Q_N(w)
       \leq\widehat R_S(g)+\lambda E.
\]
The teacher's loss lies in $[0,(U+1)^2]$ and has variance at most
$(U+1)^2R(g)$.  Bernstein's inequality therefore gives
$\widehat R_S(g)\leq R(g)+b_n(g,\delta)$ with probability at least
$1-\delta/2$.  Intersecting the two events proves the right-hand side
of~\eqref{eq:raw-oracle} for every global minimizer, at every width
$N\geq m$.

For completeness, the passage to the posterior also holds for raw
risk.  At fixed $n,N,S$, the objective is continuous and coercive.
Its minimizer set is compact, and its near-minimum sublevel sets shrink
to that set.  The risk $w\mapsto R(z_w)$ is continuous on every compact
parameter set by dominated convergence, using
$|z_w|\leq\sqrt{Q_N(w)}$.  Moreover,
\[
 R(z_w)\leq2Q_N(w)+2.
\]
To make the required tail control explicit, write
$F(w)=Q_N(w)+\lambda^{-1}\widehat R_S(z_w)$ and
$F_* = \min_w F(w)$.  For every $\eta>0$,
\begin{equation}
 \lim_{\sigma\downarrow0}
 \int (1+Q_N(w))\,
       \mathbf 1_{\{F(w)>F_*+\eta\}}\,
       \Pi_{\sigma,N,S}(dw)=0.
 \label{eq:raw-weighted-laplace-tail}
\end{equation}
Indeed, a fixed ball around a minimizer lower-bounds the normalizer by
its volume times $\exp[-(F_*+\eta/4)/(2\sigma^2)]$.
Put $\rho=\eta/[4(F_*+\eta)]$.  On the indicated tail,
$F\geq(1-\rho)(F_*+\eta)+\rho Q_N
      =F_*+3\eta/4+\rho Q_N$.
The remaining integral is a quadratic Gaussian integral of $1+Q_N$,
which grows or decays at most polynomially in $\sigma$ and is dominated
by the factor $\exp[-\eta/(4\sigma^2)]$ after division by the
normalizer.  This proves~\eqref{eq:raw-weighted-laplace-tail}.
Consequently the limiting posterior average of the risk is at most its
maximum over the minimizer set, which proves~\eqref{eq:raw-oracle}.  Jensen's inequality and Tonelli's
theorem give the conclusion for the ordinary posterior mean.  Clipping
can only decrease squared loss for labels in $[-1,1]$; applying the same
facts gives the other prediction conventions.

Finally,~\eqref{eq:raw-explicit-schedule} satisfies
$\lambda_n\geq1/n$ and the other lower bounds in
\eqref{eq:raw-calibration}.  Its leading power is $n^{-1/H}$.
The terms $r_n$ have powers at least $n^{-1/2}$, up to logarithms, and
$q_n/\lambda_n$ has power at least $n^{-1+1/H}$.  Both are no slower than
$n^{-1/H}$ when $H\geq2$.
\end{proof}

The teacher in this theorem may be selected using the data law and
$n$, but must be fixed before the sample is drawn.  In particular, it
may minimize the population risk over a specified compact class of
bounded-weight teachers.  Such a class contains the zero network, so
its minimizer has risk at most one.  This interpretation gives a sharp
oracle inequality with leading coefficient one, rather than a
realizability guarantee.  The event is simultaneous over learner widths;
no simultaneity over all teachers is needed for this interpretation.

\subsection{One hidden layer}

The same smooth-loss argument gives the rate $\widetilde O(n^{-1/2})$
for raw predictions when $H=1$.  Define $p=1/2$ and the quantities
$A_n,a_n,q_n,r_n$ by~\eqref{eq:raw-constants}.  The following explicit
version records the different absorption step.

\begin{theorem}[Raw predictions with one hidden layer]
\label{thm:raw-gaussian-one-hidden}
Let $H=1$ and suppose
\[
 1/n\leq\lambda\leq1,
 \qquad \lambda\geq2C_0(a_n^2+8q_n).
\]
For every fixed teacher $g$, with probability at least $1-\delta$,
simultaneously for every $N\geq m$,
\begin{align}
 \limsup_{\sigma\downarrow0}\int R(z_w)\,\Pi_{\sigma,N,S}(dw)
 \leq R(g)+\lambda E+r_n
       +\frac{C_0^2(a_n+2\sqrt{2q_n})^2}{2\lambda}
       +b_n(g,\delta).
 \label{eq:raw-one-hidden-oracle}
\end{align}
The ordinary posterior mean and the clipped prediction conventions
satisfy the same bound.  For example,
\[
 \lambda_n=32C_0(a_n+\sqrt{q_n}+a_n^2+q_n)
\]
satisfies the lower bounds and gives the rate
$\widetilde O(n^{-1/2})$ for a fixed teacher, whenever $n$ is sufficiently
large that $\lambda_n\leq1$.
\end{theorem}
\begin{proof}
The same simultaneous shell event yields~\eqref{eq:raw-shell-bound}.
For $p=1/2$, its right-hand side is
\[
 C_0(a_n+2\sqrt{2q_n})\sqrt k+C_0(a_n^2+8q_n)k.
\]
The second term is at most $\lambda k/2$.  Young's inequality bounds
the first by
\[
 \frac{\lambda k}2+
       \frac{C_0^2(a_n+2\sqrt{2q_n})^2}{2\lambda}.
\]
The rest is the teacher comparison, scalar concentration, and Laplace
argument from the preceding proof.  For the stated schedule,
$a_n=\widetilde O(n^{-1/2})$ and
$q_n=\widetilde O(n^{-1})$, giving the claimed rate.
\end{proof}

If a teacher-energy upper bound $E_0\geq E$ is specified in advance,
there is a useful tuning that improves the teacher dependence.  Put
\begin{gather*}
 V_n=\frac{C_0^2(a_n+2\sqrt{2q_n})^2}{2},\qquad
 v_n=\max\{2C_0(a_n^2+8q_n),\,1/n\},\\
 \lambda_n=\max\left\{\sqrt{\frac{V_n}{E_0+1}},\,v_n\right\}.
\end{gather*}
Whenever $\lambda_n\leq1$, Theorem~\ref{thm:raw-gaussian-one-hidden}
applies, and
\[
 \lambda_n E+\frac{V_n}{\lambda_n}
   \leq2\sqrt{V_n(E_0+1)}+v_nE.
\]
At fixed input dimension, this contribution is
$\widetilde O(\sqrt{(E_0+1)/n}+E_0/n)$.
For example, for population-optimal teachers whose width and weights
are both at most $n^\alpha$, one may take $E_0=O_d(n^{4\alpha})$.
The teacher risk is then at most one, and all terms in
\eqref{eq:raw-one-hidden-oracle} tend to zero for every $\alpha<1/4$,
with leading power $n^{2\alpha-1/2}$.
